\section{Estimation: staged continuation to the global fit}
\label{sec:estimation}

\subsection{The target and the three difficulties}
The scientific object is the single global posterior $p(\Theta\given X_{\mathrm{obs}})$ with
$\Theta=\{\Lambda,\alpha,\theta,\Phi,\text{latent}\}$ for all eligible dimensions, all indicators, all
likelihoods, on the full V0 sample. Fitting it cold is hard for three concrete reasons, each a failure mode:
\begin{enumerate}\setlength{\itemsep}{1pt}
  \item \textbf{Funnel geometry at scale.} $\sim$9{,}000 patients $\times$ $\sim$8 latent dimensions
  $\approx 70{,}000$ per-patient latent parameters with hierarchical variance produce Neal's funnel and NUTS
  divergences.
  \item \textbf{Weak identification.} $G$ + correlated specifics + cross-loadings is identified only up to
  rotation/sign/label; on sparse, cohort-structured missingness the posterior can be multimodal.
  \item \textbf{Mixed-likelihood stiffness.} Gaussian + ordinal + Bernoulli + negative-binomial across
  cohort-specific missingness at full $N$ is a large, stiff joint model.
\end{enumerate}

\subsection{Continuation (homotopy) staging}
Staging is \emph{not} a scientific sequence of results --- it is a continuation method. We solve an easy,
well-conditioned sub-model, then deform toward the full target one source of difficulty at a time, carrying
each converged posterior forward as the warm-start for the next. Formally we fit a nested sequence
$\mathcal{M}_1\subset\mathcal{M}_2\subset\dots\subset\mathcal{M}_5=\text{global}$ and initialise
$\theta_{\text{init}}^{(s+1)}=\E[\theta\given\mathcal{M}_s]$.

\begin{figure}[h]\centering
\resizebox{\linewidth}{!}{%
\begin{tikzpicture}[>={Stealth[length=2.6mm]},node distance=5mm,
   st/.style={draw=faceaccent,line width=0.8pt,rounded corners=2.5pt,fill=white,
     minimum height=20mm,text width=27mm,align=center,font=\sffamily\footnotesize,inner sep=3pt},
   stc/.style={draw=facekr,line width=1.0pt,rounded corners=2.5pt,fill=facekr!8,
     minimum height=20mm,text width=27mm,align=center,font=\sffamily\footnotesize,inner sep=3pt},
   stp/.style={draw=facecav,line width=1.0pt,rounded corners=2.5pt,fill=facecav!8,
     minimum height=20mm,text width=27mm,align=center,font=\sffamily\footnotesize,inner sep=3pt}]
  \node[stc] (s1) {\bfseries S1\\\scriptsize continuous core\\\scriptsize $G$+cog/met/inf/sleep\\[2pt]\scriptsize full $N$ \,\textbullet\, $\Phi=I$};
  \node[stc,right=of s1] (s2) {\bfseries S2\\\scriptsize $+\,\Phi$ + MADRS/\\\scriptsize QIDS/STAI windows\\[2pt]\scriptsize full $N$};
  \node[st,right=of s2] (s3) {\bfseries S3\\\scriptsize $+$ developmental\\\scriptsize $+$ suicidality (mixed)\\[2pt]\scriptsize subsample};
  \node[st,right=of s3] (s4) {\bfseries S4\\\scriptsize $+$ anhedonia\\\scriptsize (thin, BP/DR)\\[2pt]\scriptsize subsample};
  \node[stp,right=of s4] (s5) {\bfseries S5\\\scriptsize all 9 dimensions\\\scriptsize $+$ correlated-$G$\\[2pt]\scriptsize \bfseries certified (largest $N$)};
  \draw[->,faceaccent,line width=0.9pt] (s1)--(s2) node[midway,above,font=\scriptsize\sffamily\color{facemute}]{warm};
  \draw[->,faceaccent,line width=0.9pt] (s2)--(s3);
  \draw[->,faceaccent,line width=0.9pt] (s3)--(s4);
  \draw[->,faceaccent,line width=0.9pt] (s4)--(s5);
\end{tikzpicture}}
\caption{Staged continuation. Each stage adds exactly one source of difficulty; the converged posterior
warm-starts the next. Green $=$ certified at full $N$; amber $=$ the reported nine-dimension global fit,
documented at the largest $N$ that mixes (cross-seed Tucker $\varphi$ 0.993). \emph{Only $\mathcal{M}_5$ is
interpreted.} S3--S4 also adjudicated developmental-risk, suicidality, and anhedonia (rejected).}
\label{fig:staging}
\end{figure}

What staging buys: \emph{fault isolation} (a failure at stage $s$ implicates only what $s$ added);
\emph{warm-starts} (the hardest fit begins in a good basin); \emph{risk quarantine} (each failure mode is
introduced at a known stage with its mitigation); \emph{stability-under-elaboration} as robustness evidence;
and a natural QC/discussion gate after each step. After the seven-dimension backbone first stabilised, two
candidates that had been \emph{deferred} (mania, substance) were folded in and the \emph{joint} nine-dimension
model re-certified --- not bolted on --- so the reported map is one internally consistent object
(Chapter~\ref{sec:results}).

\begin{invariantbox}[The interpretation rule]
\textbf{Only $\mathcal{M}_5$ (the global fit) is interpreted or reported.} Stages S1--S4 are
convergence/identification checkpoints with \emph{provisional} reads. Publishing an intermediate stage is
exactly the error to avoid --- a prior iteration concluded ``no general factor'' from a sub-model that omitted
the severity anchors. Each stage emits a report and is followed by a discussion gate before advancing.
\end{invariantbox}

A stage advances only through the acceptance gate: $\widehat{R}\le 1.01$; $\mathrm{ESS}\ge 400$ (bulk and
tail); $0$ divergences; no Heywood case ($\max_j|\lambda_j|$ within a cap); and posterior-predictive checks not
grossly violated.

\subsection{The Gaussian-block marginal, and its identity with FIML}
For continuous indicators the latent factors can be integrated out (\emph{marginalized}) analytically ---
averaged over rather than estimated one patient at a time, which is what removes the per-patient ``funnel''
that otherwise makes these models hard to sample. Stacking the continuous
loadings into $\Lambda$, the full factor covariance into $\Phi$, and the residual variances into
$\Psi=\diagop(\sigma_j^2)$, the marginal over patient $i$'s observed continuous sub-vector is Gaussian with
\begin{equation}
\label{eq:marginal}
X_{i,\Obs_i}\sim\Normal\!\left(\mu_{i,\Obs_i},\;\Sigma_{\Obs_i\Obs_i}\right),
\qquad \Sigma=\Lambda\,\Phi\,\Lambda\Tn+\Psi,\quad \mu_i=\alpha+B c_i .
\end{equation}
Summing $\log\Normal(\cdot)$ over each patient's observed sub-vector is \emph{exactly the FIML objective}.

\begin{methodsbox}[Why this makes a separate FIML confirmation redundant]
The marginalized Bayesian model and a full-information maximum-likelihood SEM optimize the \emph{same}
observed-data likelihood; they differ only by the presence of priors. A flat-prior MAP therefore equals the
MLE equals the FIML solution. Hence (i) FIML adds no independent estimator evidence beyond a prior-free refit
(Chapter~\ref{sec:validation}), and (ii) this marginal is also the computational route that removes the
per-patient latent funnel for the Gaussian block while keeping the full sample.
\end{methodsbox}

\subsection{The low-rank Woodbury likelihood}
Evaluating Eq.~\eqref{eq:marginal} naively costs a determinant and a solve on $\Sigma_{\Obs_i\Obs_i}$, which
lives in indicator space (up to $71\times 71$ for a well-measured patient) --- and NUTS needs this, with its
gradient, many times per draw. But $\Sigma=\Psi+\Lambda\Lambda\Tn$ has a diagonal-plus-low-rank structure
($\Psi$ diagonal; $\Lambda\Lambda\Tn$ rank $F=K{+}1$) --- per-indicator noise plus a handful of shared factor
directions. The \textbf{Woodbury identity} (with the matrix-determinant lemma) is the standard algebraic
shortcut for exactly this shape: it lets the costly determinant and solve be carried out in the small
\emph{factor} space rather than the full indicator space --- the work moves into a room a few-by-a-few wide
instead of seventy-by-seventy. With $W=\Psi^{-1}$ (the observed-cell precision; $W_{ij}=0$
where indicator $j$ is missing for patient $i$),
\begin{equation}
A_i \;=\; I_F + \Lambda_{\Obs_i}\Tn\,W_i\,\Lambda_{\Obs_i}\in\R^{F\times F},
\qquad b_i=\Lambda_{\Obs_i}\Tn W_i x_{i,\Obs_i},
\end{equation}
\begin{equation}
\log|\Sigma_i| = \log|\Psi_i| + \log|A_i|,
\qquad
x\Tn\Sigma_i^{-1}x = x\Tn W_i x - b_i\Tn A_i^{-1} b_i,
\end{equation}
so each patient's marginal log-likelihood is
\begin{equation}
\label{eq:woodbury}
\ell_i = -\tfrac{1}{2}\Big[\, k_i\log 2\pi + \log|\Psi_i| + \log|A_i| + x_i\Tn W_i x_i - b_i\Tn A_i^{-1} b_i \,\Big],
\quad k_i=|\Obs_i| .
\end{equation}
The only matrix factorized is $A_i$, of size $F\times F$ (e.g.\ $5\times 5$ at S1) rather than up to
$68\times 68$. Correlated factors fold in through the reparameterization $\widetilde\Lambda=\Lambda\,\mathrm{chol}(\Phi)$,
so $\Lambda\Phi\Lambda\Tn+\Psi=\widetilde\Lambda\widetilde\Lambda\Tn+\Psi$ and the same kernel runs unchanged.
The likelihood is added to the model as a single potential $\sum_i\ell_i$; NUTS then samples only the
loading/residual parameters --- the $\sim$45{,}000 patient latents are gone.

\begin{figure}[h]\centering
\woodburydiagram
\caption{The mechanism behind full-sample tractability. \emph{Left:} the explicit model instantiates a latent
vector for every patient ($\approx$70{,}000 parameters) and evaluates the Gaussian likelihood with a
determinant and solve on the observed sub-covariance $\Sigma_{\mathrm{obs}}$ --- a dense matrix in
\emph{indicator} space, up to $71\times 71$. \emph{Right:} integrating the Gaussian latents out and applying
the Woodbury identity moves the only matrix that must be factorized, $A_i$, into \emph{factor} space
($\approx 5\times 5$), and removes the per-patient funnel entirely. Same likelihood, vastly cheaper geometry.}
\label{fig:woodbury}
\end{figure}

\begin{keyresult}[Full-sample inference without imputation is affordable]
The marginalized Woodbury parameterization fits the S1 continuous core on \emph{all} $N=9{,}013$ patients
($J=68$, $415{,}531$ observed cells), with no imputation and no subsampling, passing every convergence gate:
$\max\widehat{R}=1.010$, $\min\mathrm{ESS}=1{,}939$, $0$ divergences. The point is not an absolute runtime but
that the full-sample, no-imputation invariant is reached by a sequence of \emph{exact} reparameterizations
(\S\ref{sub:cost}) rather than by approximation or by selecting an easier sub-population --- a cold
explicit-latent fit of the same model does not converge at this scale.
\end{keyresult}

\paragraph{Recovering patient scores.} Marginalization does not discard the coordinates: for a fixed posterior
draw the conditional $z_i\given x_{i,\Obs_i}$ is Gaussian with precision $A_i$ and mean $A_i^{-1}b_i$ --- the
\emph{same} $A_i,b_i$ used in Eq.~\eqref{eq:woodbury}. Fitting and scoring thus separate cleanly: fit the map
once on the full sample, then project any patient's observed cells onto it (with uncertainty), so a downstream
cohort can be scored without driving the measurement fit.

\subsection{The mixed-likelihood frontier (S3--S5)}
Binary/ordinal/count indicators cannot be Gaussian-marginalized. The engine uses a \emph{partial collapse}: the
factors touching non-Gaussian indicators ($f_e=\{G,\text{suicidality},\text{developmental}\}$) carry explicit
non-centered latents, while the pure-continuous specifics ($f_m=\{$cognition, metabolic, inflammatory, sleep$\}$)
stay marginalized and are coupled to $f_e$ through the shared $\Phi$ by the Gaussian conditional
\begin{equation}
f_m\given f_e\sim\Normal(M f_e,\,S),\qquad M=\Phi_{me}\Phi_{ee}^{-1},\quad S=\Phi_{mm}-\Phi_{me}\Phi_{ee}^{-1}\Phi_{em}.
\end{equation}
The continuous block becomes a per-patient-mean Woodbury (mean $B f_e$, residual covariance
$\Lambda_m S\Lambda_m\Tn+\Psi$); the non-Gaussian indicators get Bernoulli/negative-binomial/ordered-logistic
likelihoods on $f_e$. This is exactly the ``partial collapse'' the methods plan flagged as the hardest step,
and it works: the mixed fit attains $0$ divergences.

\begin{methodsbox}[A reparameterization ladder, not hardware, resolved the mixing frontier]
With $0$ divergences but slow mixing, the mixed-likelihood stages were weak-identification \emph{ridges}, not
funnels --- so the fix is reparameterization, not more compute. Two moves did it: (i) tightening every explicit
specific's $\to G$ cross-loading toward zero (the dimensions that touch non-Gaussian items are $\approx\perp G$
anyway, so this removes a ridge without touching the biology$\perp G$ estimand), and (ii) building $\Phi$ from
a \emph{unit-row} Cholesky factor so it stays a valid correlation matrix (the same fix as the $\Phi=LL\Tn$ bug,
Chapter~\ref{sec:engineering}). The explicit-latent block remains the documented mixing limit, so the joint map
is documented at the \emph{largest $N$ that mixes} with a cross-seed resample-stability guard.
\end{methodsbox}

\subsection{What makes the fit expensive --- and the optimizations that remove it}
\label{sub:cost}
It is worth being explicit about where the time goes. The cost of one NUTS draw factorizes as
\[
\text{cost per draw}\;\sim\;\underbrace{\big(\text{\# sampled parameters}\big)}_{\text{geometry}}\;\times\;
\underbrace{\big(\text{leapfrog steps per draw}\big)}_{\text{step size}}\;\times\;
\underbrace{\big(\text{cost of one gradient}\big)}_{\text{linear algebra}},
\]
and a cold, explicit-latent fit is large on all three. The \emph{geometry} factor is enormous: instantiating
$G_i$ and $D_{ik}$ for every patient adds $\sim$45{,}000--70{,}000 latent parameters, and their hierarchical
variance creates Neal's funnel, which forces tiny steps and produces divergences (so the \emph{step-size}
factor is bad too). The \emph{linear-algebra} factor is also large: every gradient evaluates the observed-data
Gaussian likelihood, which naively requires a determinant and a solve on each patient's observed sub-covariance
$\Sigma_{\Obs_i\Obs_i}$ --- a dense matrix in indicator space, up to $71\times 71$, across all $\sim$9{,}000
patients --- and NUTS performs this many times per retained draw, so the three factors multiply.

The engine attacks each factor with a transformation that leaves the inferential target unchanged
(Table~\ref{tab:speedups}); the one approximation in the stack is confined to the heavier checkpoint stages and
is labelled as such.

\begin{table}[h]\centering\small\sffamily
\caption{The optimization stack. Every row preserves the inferential target (the same posterior, or the same
observed-data likelihood) except the last --- an explicitly labelled approximation used only where exact
marginalization is unavailable.}
\label{tab:speedups}
\begin{tabular}{L{2.7cm} L{6.1cm} L{5.0cm}}
\toprule
\hdr{factor attacked} & \hdr{optimization} & \hdr{effect}\\
\midrule
\rowcolor{facebg} geometry & marginalize the Gaussian latents (Eq.~\ref{eq:marginal}) & removes $\sim$45--70k parameters; eliminates the funnel; the MCMC state drops to a few hundred structural parameters\\
linear algebra & Woodbury identity $+$ determinant lemma (Eq.~\ref{eq:woodbury}) & the per-patient solve moves from indicator space ($\le 71^2$) to factor space ($F^2$, e.g.\ $5^2$)\\
\rowcolor{facebg} linear algebra & pattern-grouped, batched assembly & $A_i$ as one matrix multiply; its Cholesky computed once per \emph{unique missingness pattern} ($\approx 2.75\times$, log-likelihood-identical)\\
step size & tuning matched to the (now benign) geometry & lower target acceptance ($0.85$--$0.90$) $+$ tree-depth cap $8$ bound the deep warmup trajectories ($\approx 2.7\times$)\\
\rowcolor{facebg} warmup length & staged warm-starts (continuation, Chapter~\ref{sec:estimation}) & each stage starts in the previous basin; shorter tuning, fewer failed trajectories\\
implementation & vectorization $+$ cached preprocessing & the $F\times F$ work is batched over patients; harmonization is not recomputed per fit\\
\midrule
\rowcolor{facecav!9} \emph{checkpoints only} & balanced random sub-sample (\emph{labelled approximation}) & the mixed stages, whose non-Gaussian block has no closed-form marginal, run at $N=4$--$5$k\\
\bottomrule
\end{tabular}
\end{table}

Two entries deserve a note. First, the tuning is counter-intuitive: because marginalization removes the funnel,
the posterior is benign (its typical set is only $\sim$21 leapfrog steps deep), so the right move is to
\emph{lower} the target acceptance and cap the tree depth --- bounding the expensive deep warmup trajectories
--- rather than the usual reflex of \emph{raising} target acceptance to chase divergences, which here would only
shrink the step size and slow a posterior that has no pathology to avoid. Second, the model is written in PyMC
for a readable specification and sampled through the NumPyro/JAX backend, which compiles and batches the
$F\times F$ linear algebra across patients. Together these turn an intractable cold fit into one that converges
at full $N$, within a standard workstation's memory, in about an hour.

\begin{caveatbox}[Why the mixed-likelihood stages cost more, and the sub-sample fallback]
The whole stack hinges on marginalizing the Gaussian latents. The non-Gaussian indicators
(binary/ordinal/count) have \emph{no closed-form latent integral}, so the suicidality and developmental factors
must keep explicit per-patient latents --- which restores the geometry cost the continuous stages had removed.
The S3--S5 checkpoints therefore run on a \emph{random, cohort-balanced} sub-sample ($N=4{,}000$--$5{,}000$;
realistic missingness, with a resample-stability check) --- never completeness- or attrition-selected, so the
missingness structure is preserved. This is the one approximation in the stack, and it is labelled as such: the
reported map targets the full sample or the largest $N$ that certifies, and a full-$N$ certified S5 is a matter
of additional sampling effort, not a change of method.
\end{caveatbox}
